% Generated by scripts/assemble.py; corrections belong in source/corrections.json.
\examhead{64}{西南大学 2024 年数学分析试卷}{64}
% p065-b16; source PDF page 65
\par\Needspace{5\baselineskip}
\blockanchor{p065-b16}
\majorquestion{一}叙述题. 每题 5 分, 共 15 分.\par
% p065-b17; source PDF page 65
\blockanchor{p065-b17}
\question{1}函数极限的柯西收敛准则.\par
% p065-b18; source PDF page 65
\blockanchor{p065-b18}
\question{2}函数 $f(x)$ 在区间 $I$ 上不一致连续的定义.\par
% p065-b19; source PDF page 65
\blockanchor{p065-b19}
\question{3}函数极限 $\lim_{x\to +\infty}f(x)$ 的归结原则.\par
% p065-b20; source PDF page 65
\par\Needspace{5\baselineskip}
\blockanchor{p065-b20}
\majorquestion{二}计算题. 每题 5 分, 共 15 分.\par
% p065-b21; source PDF page 65
\blockanchor{p065-b21}
\question{1}求不定积分 $\int \sin (\ln x) \, \mathrm{d}x$.\par
% p065-b22; source PDF page 65
\blockanchor{p065-b22}
\question{2}求极限 $\lim_{x\to +\infty}\left[\frac{(1 + x)^{\frac{1}{x}}}{e}\right]^{\frac{1}{x}}.$\par
% p065-b23; source PDF page 65
\blockanchor{p065-b23}
\question{3}设 $f$ 具有二阶连续的偏导数, 且 $z = f(x^{2}y^{2}, xy)$, 求 $z_{xy}$.\par
% p065-b24; source PDF page 65
\par\Needspace{5\baselineskip}
\blockanchor{p065-b24}
\majorquestion{三}(15 分) 设有界函数 $f(x)$ 在 $[a, b]$ 上的不连续点的全体为 $\{x_n \mid n = 1, 2, \cdots\}$, 且 $\lim_{n \to \infty} x_n = a$, 证明: $f(x)$ 在 $[a, b]$ 上可积.\par
% p066-b00; source PDF page 66
\par\Needspace{5\baselineskip}
\blockanchor{p066-b00}
\majorquestion{四}(15 分) 设数列 $\{a_{n}\}$ 严格递减, 且 $a_{n} > 0$. 证明: $\lim_{n \to \infty} a_{n} = 0$ 的充要条件是级数 $\sum_{n=1}^{\infty} \left(1 - \frac{a_{n+1}}{a_{n}}\right)$ 发散.\par
% p066-b01; source PDF page 66
\par\Needspace{5\baselineskip}
\blockanchor{p066-b01}
\majorquestion{五}(15 分) 已知 $f(x)$ 在 $(- \infty, +\infty)$ 上二阶可导, 且 $f''(x) \leq 0$, $\lim_{x \to 0} \frac{f(x)}{2x} = 2$, 证明: $f(x) \leq 4x$.\par
% p066-b02; source PDF page 66
\par\Needspace{5\baselineskip}
\blockanchor{p066-b02}
\majorquestion{六}(15 分) 用区间套定理证明确界原理.\par
% p066-b03; source PDF page 66
\par\Needspace{5\baselineskip}
\par\Needspace{6\baselineskip}
\blockanchor{p066-b03}
\majorquestion{七}(15 分) 设函数 $f(x)$ 在 $[1, +\infty)$ 上连续, 恒正且单调递减, 记\par
\nopagebreak[4]
% p066-b04; source PDF page 66
\blockanchor{p066-b04}
\displayformula{u _ {n} = \sum_ {k = 1} ^ {n} f (k) - \int_ {1} ^ {n} f (x) \mathrm{d} x.}
% p066-b05; source PDF page 66
\blockanchor{p066-b05}
证明数列 $\{u_n\}$ 收敛.\par
% p066-b06; source PDF page 66
\par\Needspace{5\baselineskip}
\blockanchor{p066-b06}
\majorquestion{八}(15 分) 设 $f(x)$ 在 $[a, +\infty)$ 上一致连续, $\varphi(x)$ 在 $[a, +\infty)$ 上连续, 且 $\lim_{x \to +\infty} [f(x) - \varphi(x)] = 0$, 证明: $\varphi(x)$ 在 $[a, +\infty)$ 上一致连续.\par
% p066-b07; source PDF page 66
\par\Needspace{5\baselineskip}
\blockanchor{p066-b07}
\majorquestion{九}(15 分) 证明半径为 $R$ 的球的体积为 $\frac{4\pi R^3}{3}$.\par
% p066-b08; source PDF page 66
\par\Needspace{5\baselineskip}
\par\Needspace{6\baselineskip}
\blockanchor{p066-b08}
\majorquestion{十}(15 分) 计算曲线积分\par
\nopagebreak[4]
% p066-b09; source PDF page 66
\blockanchor{p066-b09}
\displayformula{\int_ {L} y \mathrm{d} x + z \mathrm{d} y + x \mathrm{d} z.}
% p066-b10; source PDF page 66
\blockanchor{p066-b10}
其中 L 是  $x^{2} + y^{2} + z^{2} = 1$  和  $x + y + z = 1$  的交线, 从 x 轴正向看去取逆时针方向.\par
